\section{Surface flux}
{\bf \Large 
\begin{tabular}{ccc}
\hline
  Corresponding author & : & Yuta Kawai\\
\hline
\end{tabular}
}

Surface fluxes represent the exchange of momentum, heat, and moisture between the surface and the lowest model level of the atmosphere.
SCALE-DG provides two surface-flux parameterizations: a constant-flux and a simple bulk-flux scheme.
In the DG framework, the surface fluxes are imposed as lower-boundary conditions at the model bottom. 

\subsection{Constant flux}
%%
This option prescribes constant surface heat and moisture fluxes independently of the atmospheric state.
On the other hand, the surface momentum flux is calculated 
using either a prescribed bulk coefficient for momentum $C_m$ or a prescribed friction velocity $u_*$. 
The surface momentum flux is writen as
%%
\begin{align}
  F_{\rho \bm{u},{\rm sfc}} = -C_m |\tilde{\bm{u}}| \widetilde{\rho \bm{u}}
\end{align}
%%
where the tilde denotes the value at the lowest model level of the atmosphere. 
When the friction velocity is prescribed, $C_m$ is diagnosed as
%%
\begin{align}
  C_m = \left(\dfrac{u_*}{|\tilde{\bm{u}}|}\right)^2.
\end{align}
%%


In SCALE-DG, the constant-flux scheme is calculated using the SCALE library module {\tt scale\_atmos\_phy\_sf\_const.F90}. 
The minimum values of the wind speed and $C_m$ can be modified from the default values (i.e., zero) to avoid numerical instability.


\subsection{Simple bulk flux scheme}
%%
This option calculates the surface fluxes 
based on constant values of the bulk drag coefficient $C_m$, the heat-transfer coefficient $C_h$, and the moisture-transfer coefficient $C_e$.
The surface momentum, sensible heat, and latent heat fluxes are calculated as
%%
\begin{align}
  F_{\rho \bm{u},{\rm sfc}} &= -C_m |\tilde{\bm{u}}| \widetilde{\rho \bm{u}},\\
  F_{{\rm SH, sfc}} &= C_h |\tilde{\bm{u}}| \tilde{\rho} (T_{\rm sfc} - \tilde{T}),\\
  F_{{\rm LH, sfc}} &= L C_e |\tilde{\bm{u}}| \tilde{\rho} (q_{\rm sfc} - \tilde{q}).
\end{align}
%%
where the tildes denote the values at the lowest model level of the atmosphere, 
and $T_{\rm sfc}$ and $q_{\rm sfc}$ are the surface temperature and surface specific humidity, respectively.
The surface specific humidity is calculated from the surface temperature and the surface pressure using the Clausius-Clapeyron equation as
%%
\begin{align}
  q_{\rm sfc} = \dfrac{\epsilon p_{\rm sat,sfc}}{p_{\rm sfc} - (1-\epsilon) p_{\rm sat,sfc}},
\end{align}
%%
where $\epsilon$ is the ratio of the molecular weight of water vapor to that of dry air,
and $p_{\rm sat,sfc}$ is the saturation vapor pressure at $T_{\rm sfc}$.

The minimum values of the wind speed can be modified from the default values (i.e., zero) to avoid numerical instability.
